\FloatBarrier

\section{Real data: CelebA and CIFAR-10 latent diffusion}
\label{sec:real-data}

\begin{figure*}[!t]
  \centering
  \begin{subfigure}{0.48\textwidth}
    \centering
    \includegraphics[width=\linewidth]{figures/celeba_final_mem_vs_d.pdf}
    \caption{CelebA final memorization.}
  \end{subfigure}\hfill
  \begin{subfigure}{0.48\textwidth}
    \centering
    \includegraphics[width=\linewidth]{figures/celeba_final_fid_vs_d.pdf}
    \caption{CelebA final FID.}
  \end{subfigure}

  \begin{subfigure}{0.48\textwidth}
    \centering
    \includegraphics[width=\linewidth]{figures/cifar10_final_mem_vs_d.pdf}
    \caption{CIFAR-10 final memorization.}
  \end{subfigure}\hfill
  \begin{subfigure}{0.48\textwidth}
    \centering
    \includegraphics[width=\linewidth]{figures/cifar10_final_fid_vs_d.pdf}
    \caption{CIFAR-10 final FID.}
  \end{subfigure}
  \caption{\textbf{Real-data latent width produces a clean quality--memory tradeoff.} Five-seed means with 95\% confidence intervals after 5M MLP diffusion steps on a diverse 1k subset. Memorization is only detectable once the VAE preserves pixel-space sample identity; past that identity threshold, increasing $\dlat$ moves the model into the excess-latent regime and reduces memorization. CelebA peaks near $\dlat=70$--$80$, while CIFAR-10 peaks later and remains high through roughly $\dlat=100$--$160$ before falling at larger widths. FID follows a separate quality tradeoff: it improves while the VAE bottleneck is being relieved, then worsens once the latent score problem becomes too wide.}
  \label{fig:real-final-mem-fid}
\end{figure*}

\subsection{VAE and diffusion training}
\label{sec:vae-training}

Real image latents introduce a second effect absent from the synthetic construction: if the VAE bottleneck is too narrow, the training images cannot be reconstructed faithfully, so pixel-space memorization is hard to detect even if the latent diffusion model fits its training codes. In the main text we therefore show compact memorization/FID summaries for both CelebA and CIFAR-10, while leaving the full per-dimension diagnostics in the appendix. The VAE validation metrics, reconstruction contact sheets, optimization losses, and full memorization/FID sweeps are in Appendix~\ref{app:vae-clean} and Appendix~\ref{app:real-diffusion-clean}: see Figures~\ref{fig:app-vae-metrics}--\ref{fig:app-vae-quality} for the VAE checks and Figures~\ref{fig:app-real-optimization}-- \ref{fig:app-real-individual-fid-cifar} for the diffusion checks. The VAE reconstruction-memorization diagnostic in Appendix Figure~\ref{fig:app-vae-metrics} gives an empirical proxy for the effective intrinsic-dimensionality threshold. We define $\widehat d_{\rm ID}^{\rm VAE}(\rho)$ as the smallest latent dimension whose immediate VAE encode/decode reconstructions are counted as memorized for at least a fraction $\rho$ of the same diverse 1k subset, using the pixel-space nearest-neighbor ratio test. In the appendix we use $\rho=0.95$ as a practical identity-preservation threshold, giving $\widehat d_{\rm ID}^{\rm VAE}(0.95) \approx 70$ for CelebA and $\approx 140$ for CIFAR-10. Below this threshold, the VAE itself averages away sample identity, so pixel-space memorization and the latent-spectrum predictor are not expected to follow the excess-buffer story cleanly.

For each dataset and each $\dlat$, we train a convolutional VAE on the full available training split rather than on the 1k diffusion subset. The encoder is a five-stage ResNet-style convolutional network with channel widths $(32,64,128,256,512)$, stride-two downsampling, GroupNorm/SiLU residual blocks, and separate linear heads for $\mu(x)$ and $\log\sigma^2(x)$. The decoder mirrors this architecture with transposed-convolution upsampling and a Tanh output. Images are represented at $32\times 32$ resolution. The VAE objective is a reconstruction term plus a KL penalty to the standard normal prior, with the final runs using the capacity-style loss $\mathcal L_{\rm VAE}=\mathcal L_{\rm recon} + \gamma |\mathrm{KL}(q(z|x)\|p(z))-C(t)|$ when enabled. The capacity target is warmed up during training, learning rate is reduced on validation plateaus, and the checkpoint with the best validation objective is used downstream.

We then freeze each VAE and train an MLP score model directly in its latent space. The diffusion subset is a farthest-first diverse set of 1k training images, encoded once by the VAE mean. The score network is a fixed-width, depth-three MLP with sinusoidal time features, hidden width 256, and output dimension $\dlat$. Keeping the hidden width fixed makes parameter count and per-step FLOPs roughly constant across the sweep, so the high-dimensional decline is not explained by simply giving wider latent spaces a wider network. Each run uses SGD with learning rate $10^{-3}$ and momentum $0.80$, trains for 5M steps over five seeds, logs optimization loss every 10k steps, and evaluates memorization and FID every 100k steps using generated samples. For the real-data memorization metric, generated latents are decoded back to images and nearest-neighbor ratios are computed in pixel space against the fixed 1k training subset, not in VAE latent space. The aggregate real-data trajectories are shown in Appendix Figure~\ref{fig:app-real-trajectories}, and the per-$\dlat$ memorization/FID panels are in Figures~\ref{fig:app-real-individual-mem}-- \ref{fig:app-real-individual-fid-cifar}. As an architecture check, Appendix~\ref{app:spatial-dit} repeats the CelebA latent-width question with a spatial VAE and DiT trunk; we keep it as an appendix pilot because changing a spatial channel count also changes the number of latent grid coordinates and their covariance structure.

\subsection{Real-data memorization and FID}
\label{sec:real-mem-fid}

On CelebA, final memorization is largest near the VAE identity threshold, around $\dlat=70$--$80$, and then falls by roughly half by $\dlat=200$. The low-$\dlat$ side is a representation-capacity regime where decoded samples do not preserve enough identity for pixel-space KNN recall to be measured reliably; the high-$\dlat$ side is the excess-latent regime, where $q$ falls from $2.9$ to $1.3$; the $q$-mechanism accounts for $1.7$--$1.9\times$ of the observed $2.1$--$3.9\times$ slowdown of the $5\%$ hitting time and for $0.28$ of the log-decline of the CelebA final fraction, and a $q$-only reading is falsified across datasets (CelebA $0.45$ vs CIFAR-10 $0.21$ final fraction at matched $q\approx1.25$). CIFAR-10 shows the same qualitative behavior with a larger identity-preservation threshold ($\widehat d_{\rm ID}^{\rm VAE}(0.95)\approx 140$): after the VAE begins preserving identity, memorization declines for the widest latents. FID has its own dimension dependence: it improves rapidly as the VAE bottleneck is relieved, is best at intermediate latent widths, and worsens again at large $\dlat$, giving a quality tradeoff rather than a monotone benefit from larger codes. The aggregate over-training memorization/FID trajectories and the full per-dimension real-data panels are shown in Appendix Figures~\ref{fig:app-real-trajectories}-- \ref{fig:app-real-individual-fid-cifar}.
